Machine learning has fundamentally transformed decision-making by
allowing algorithms to discover intricate and
meaningful patterns directly from data, without relying explicitly on
predefined rules. Rather than encoding
human expertise through manual programming, machine learning
algorithms generalize from examples. This inductive
process, which generalizes observed cases to unseen scenarios,
enables the algorithm to identify underlying patterns
from historical examples and predict future outcomes. However, such a
reliance on historical data inherently carries
risks, especially when data reflects existing societal biases,
stereotypes, or inequalities.

\subsection{Algorithmic Bias}\label{subsec:algorithmic_bias}

Algorithmic bias refers to systematic and repeatable errors or unfair
outcomes produced by machine learning models
due to biases embedded within the training data or algorithm design.
According to existing literature
\cite{barocas2016big,pessach2020algorithmic}, algorithmic bias
commonly originates from:

\begin{itemize}
  \item \textbf{Biases inherent in training datasets:} These biases
    result from biased human decisions, measurement errors,
    reporting inaccuracies, or historical prejudices embedded in
    datasets. Machine learning algorithms, aiming at optimizing
    prediction accuracy, often replicate these biases.
  \item \textbf{Biases due to missing data:} When datasets lack
    sufficient representation from certain groups or have
    significant data omissions, the resulting models fail to
    accurately represent the target population.
  \item \textbf{Algorithmic optimization bias:} Typical optimization
    objectives, such as minimizing aggregate prediction
    errors, tend to favor majority groups, often leading to poorer
    performance for minority groups.
  \item \textbf{Bias from proxy attributes:} Non-sensitive attributes
    may indirectly capture sensitive information
    (e.g., race, gender, age), unintentionally introducing biases
    even when sensitive attributes are explicitly excluded
    from the dataset.
\end{itemize}

This means that algorithmic bias is a fundamental and unavoidable
fact when developing machine learning algorithms. However, not all is
lost, by carefully defining what it is we mean when we require an
algorithm to be "fair", that is without bias, machine learning
research has made immense strides in making these algorithms more
accurate, correct, and true representations of real-world scenarios.
\subsection{Defining Fairness in Machine
Learning}\label{subsec:fairness_definitions}

Given the increasing use of machine learning in high-stakes domains,
rigorous fairness definitions have emerged to guide
algorithmic development. These definitions typically fall into two
broad categories: individual fairness and group fairness.

\paragraph{Individual Fairness.}
Individual fairness requires models to produce similar outputs for
similar individuals, where similarity is assessed based
on relevant non-sensitive features. Formally, individual fairness can
be articulated using Lipschitz continuity
as follows \cite{dwork2012fairness}:
\[
  d(\text{output}(x), \text{output}(x')) \leq \rho \cdot d(x, x'),
\]
where \(x\) and \(x'\) are individuals with comparable non-sensitive
attributes, and \(\rho\) is a small constant. This definition
emphasizes the fair treatment of similar cases on an individual
basis, ensuring minimal unjustified variability. Heuristically, we
find individual fairness to be a stronger criterion for fairness, but
in practice, it is much harder to implement.

\paragraph{Group Fairness.}
Group fairness demands that statistical outcomes of algorithms be
equitable across predefined demographic groups.
This approach explicitly acknowledges and attempts to rectify
societal disparities. Notable metrics
for group fairness include:

\begin{itemize}
  \item \textbf{Demographic Parity:} Ensures equal rates of positive
    predictions across demographic groups:
    \[
      P(R = 1 \mid A = a) = P(R = 1 \mid A = b),
    \]
    where \(R\) denotes the model's prediction and \(A\) represents a
    protected attribute (e.g., race, gender).

  \item \textbf{Equal Opportunity:} Requires equal true positive
    rates among groups, ensuring fairness in the
    allocation of positive outcomes given actual positives:
    \[
      P(R = 1 \mid Y = 1, A = a) = P(R = 1 \mid Y = 1, A = b),
    \]
    with \(Y\) representing the true outcome.

  \item \textbf{Equalized Odds:} Further requires equal true positive
    and false positive rates across groups,
    encompassing both success and error equity:
    \[
      P(R = 1 \mid Y = y, A = a) = P(R = 1 \mid Y = y, A = b), \quad
      \forall y \in \{0, 1\}.
    \]
\end{itemize}

Each of these metrics presents trade-offs, and no universal criterion
exists to satisfy all simultaneously,
leading to a fundamental tension explored later.

\subsection{Real-World Instances and Ethical Dimensions of
Algorithmic Bias}\label{subsec:real_world_bias}

Concrete examples vividly illustrate the risks associated with biased
algorithms. One widely cited example is the COMPAS algorithm,
frequently used in criminal justice for predicting recidivism.
Investigations revealed significant racial biases, incorrectly
labeling Black defendants as high-risk at nearly twice the rate of
White defendants who later re-offended ~[\cite{angwin2016machine}].
Similarly, facial recognition software has consistently demonstrated
higher error rates for darker-skinned individuals, exacerbating risks
of racial profiling and wrongful identification.

Algorithmic biases also permeate employment contexts, where
historical data reflecting past hiring decisions embed biases against
women or minority groups, perpetuating discrimination through
ostensibly neutral automated decision-making systems. When training
algorithms on historical employment records, biases and stereotypes
embedded in the data disproportionately disadvantage female candidates.

A particularly telling example emerges from machine translation.
Consider translating sentences from English to Turkish and then back
into English, as illustrated in Figure~\ref{fig:eng-to-turkish}.
Turkish pronouns are gender-neutral,
but when translated back into English, gender-specific pronouns are
inferred based on statistical associations. As a result, occupations
stereotypically associated with men—such as ``engineer'' or
``doctor''—are translated back using male pronouns, while occupations
stereotypically associated with women—such as ``nurse''—return female
pronouns. This phenomenon arises from two biases embedded in training
datasets: first, real-world labor market statistics reflecting historical
occupational distributions and, second, the ``male-as-norm'' bias, whereby
male pronouns are preferentially selected when gender is ambiguous or
unknown~\cite{caliskan2017semantics}.

\begin{figure}[h]
  \centering
  \includegraphics[width=0.9\textwidth]{sections/background/eng-to-turkish.png}
  \includegraphics[width=0.9\textwidth]{sections/background/turkish-to-eng.png}
  \caption{Google translate screenshot of gender-specific English sentences into
  gender-neutral Turkish and then back to English.}
  \label{fig:eng-to-turkish}
\end{figure}

Attempts to mitigate biases by removing explicitly sensitive
attributes (such as gender or race) from training datasets frequently
fall short due to \textit{proxy variables}. Proxy attributes, for example,
the age at which individuals start programming, can inadvertently
encode sensitive information like gender, reinforcing biases even in
their absence. Additionally, biases due to disparities in sample
sizes among demographic groups lead to poorer model performance for
minority groups, reinforcing systematic inequities~\cite{barocas2016big}.

Beyond technical considerations, fairness intersects profoundly with
ethical and legal imperatives. Legal frameworks such as Title VII of
the U.S. Civil Rights Act mandate nondiscrimination in employment.
Similarly, regulatory initiatives like the European Union’s GDPR and
proposed AI Act embed fairness, transparency, and equity into
regulatory requirements for algorithmic systems.

Ethically, ensuring fairness aligns closely with broader principles
of justice and equity, particularly as algorithmic systems
increasingly influence societal outcomes such as employment,
education, and criminal justice. Effective fairness interventions
help prevent reinforcing historical injustices, foster social equity,
and maintain public trust in algorithm-driven decision-making.

\subsection{Challenges, Limitations, and Future
Directions}\label{subsec:challenges_future}
(TODO: citations here need to be filled in carefully I relied a bit
on yt tAjFuhkiV2c)

The growing emphasis on fairness in algorithmic decision-making has
significantly advanced our understanding of bias mitigation.
Nevertheless, this field continues to face considerable challenges
and inherent limitations.

One prominent difficulty arises from conflicting fairness criteria.
It is mathematically impossible (TODO: why? should I explain here?)
to simultaneously satisfy all fairness definitions — like
Demographic Parity, Equal Opportunity,
and Equalized Odds — highlighting the necessity for context-specific
fairness solutions. For example, attempts to apply broad fairness
concepts, initially developed within supervised learning frameworks,
to clustering tasks often encounter mismatches in meaning. [cite] Individual
fairness definitions emphasizing distributional equity may not
translate effectively into clustering contexts, where groups or
clusters often lack inherent meaning until assigned. This mismatch
underscores the need to carefully adapt fairness criteria from
supervised learning to unsupervised scenarios. [cite]

Moreover, algorithmic fairness interventions rarely function in
isolation. They typically constitute components within broader
socio-technical systems, necessitating careful consideration of both
upstream inputs and downstream impacts. The removal of sensitive
variables, a common fairness strategy, can inadvertently lead to
unintended consequences. For instance, the practice of "Ban the Box,"
intended to eliminate employment discrimination by prohibiting
questions about criminal history, inadvertently increased racial
discrimination as employers began using race as a proxy variable. [cite]
This underscores the complexity and potential pitfalls inherent in
algorithmic fairness interventions and highlights the importance of
anticipating and managing unintended downstream consequences.

Further complications emerge from the dynamic nature of real-world
data. In applications such as school districting or political
redistricting, clustering algorithms are applied to dynamic
populations, where individuals may relocate in response to
algorithmic interventions. Historical efforts like school busing
aimed at integrating racially segregated districts illustrate how
algorithmic clustering solutions can unintentionally disrupt
communities or exacerbate segregation through mechanisms such as
"white flight." [cite] This historical context reveals the critical need to
engage deeply with domain-specific constraints, legal frameworks, and
stakeholder needs, reinforcing that algorithmic solutions must be
cognizant of broader historical and societal dynamics.

Indeed, addressing these multidimensional challenges requires
interdisciplinary collaboration beyond computer science. Integrating
insights from fields such as sociology, criminology, economics, law,
and ethics is critical. For instance, fairness research often
overlooks valuable contributions from disciplines like criminology or
education policy, which provide nuanced understandings of systemic
inequities and practical constraints. Collaboration with experts from
these domains can guide the appropriate adaptation of algorithmic
methods to complex social contexts, ensuring solutions align closely
with practical realities and ethical standards.

Another essential consideration is stakeholder engagement. Too often,
fairness solutions are developed paternalistically, without
adequately involving those directly impacted.[cite] Engaging
stakeholders—including affected communities, policy experts, and
practitioners—in defining fairness and assessing interventions can
prevent misguided assumptions and ensure that algorithmic systems
genuinely serve intended beneficiaries. This is particularly evident
in sensitive domains such as criminal justice, healthcare, and
education, where the risk of inadvertently causing harm or
perpetuating injustices remains high.

Finally, reliance on commonly used benchmark datasets, such as the
COMPAS, Adult, and German Credit datasets, introduces risks of
replicating inherent data biases and inaccuracies. Issues such as
noisy demographic labels, inappropriate or misleading features, and
the misalignment of fairness labels highlight critical weaknesses in
the existing empirical evaluation landscape. Consequently, rigorous
methodological scrutiny and the development of better benchmarks
reflecting real-world complexities and accurate demographic
information are urgently needed.

In summary, future directions in algorithmic fairness research must
navigate inherent mathematical and practical complexities through
rigorous interdisciplinary collaboration, careful stakeholder
engagement, and meticulous empirical practices. By integrating
technical algorithmic approaches with broader societal, legal, and
ethical considerations, researchers can develop more robust,
equitable, and practically viable solutions to mitigate algorithmic
biases and their profound societal implications.

(TODO: connect to the next section here? required? if so, how?)